\section{Kịch bản Kiểm thử và Đánh giá Hiệu năng}

\subsection[Tổng quan hệ thống Kiểm thử tự động]{Tổng quan hệ thống Kiểm thử tự động (Test Framework)}

\subsubsection{Cơ sở lý thuyết về phương pháp luận kiểm thử}

Trong chu trình phát triển phần mềm hiện đại, Phương pháp Phát triển Hướng Kiểm thử (Test-Driven Development - TDD) và hệ thống Tích hợp Liên tục (Continuous Integration - CI) đóng vai trò then chốt trong việc duy trì độ ổn định và chất lượng mã nguồn. Thay vì triển khai kiểm thử thủ công (Manual Testing) vốn tiêu tốn tài nguyên và dễ phát sinh sai sót, hệ thống được trang bị một hệ sinh thái kiểm thử tự động (Automated Testing Framework) tích hợp trực tiếp vào quy trình đóng gói.

Hệ thống được thiết kế bao gồm hai cấp độ kiểm thử chuyên sâu:
\begin{itemize}[label={-}]
    \item \textbf{Kiểm thử Đơn vị (Unit Testing):} Được thực thi tại mức mã nguồn (Code-level) nhằm cô lập và xác thực tính đúng đắn của từng hàm hoặc thành phần (Module) cụ thể. Kỹ thuật này cho phép phát hiện sớm các lỗ hổng logic từ giai đoạn phát triển ban đầu (Shift-left Testing).
    \item \textbf{Kiểm thử Tích hợp (Integration/E2E Testing):} Xác minh luồng giao tiếp dữ liệu liên thông giữa ứng dụng máy khách và máy chủ, mô phỏng các kịch bản tương tác thực tế dựa trên phương pháp hộp đen (Black-box Testing).
\end{itemize}

\subsubsection{Cấu trúc Test Framework trong thực tế mã nguồn}

Thay vì phụ thuộc vào các thư viện bên thứ ba đồ sộ (như Google Test hay Catch2) làm gia tăng chi phí biên dịch, dự án đã tự chủ thiết kế một vi khung kiểm thử (Micro Test Framework) bằng C++ thuần túy thông qua các tệp nền tảng \texttt{TestMacro.h} và \texttt{test\_framework.h}. Hệ thống này cung cấp các hàm Macro tiêu chuẩn hóa như \texttt{ASSERT\_TRUE}, \texttt{ASSERT\_EQ}, và \texttt{ASSERT\_THROWS} để đánh giá các biểu thức logic tĩnh.

Thông qua chỉ thị \texttt{enable\_testing()} được khai báo trong \texttt{CMakeLists.txt}, hệ thống đăng ký và vận hành tự động một tập hợp các kịch bản kiểm thử (Test Suites) độc lập:

\begin{itemize}[label={-}]
    \item \textbf{Khối Kiểm thử Mật mã (\texttt{test\_crypto}):} Đo lường độ phức tạp thời gian thực thi (Time Complexity) và xác thực tính đồng nhất của dữ liệu sau quá trình mã hóa và giải mã lai giữa phương thức bất đối xứng RSA-2048 và luồng khối AES-256.
    \item \textbf{Khối Kiểm thử Xác thực (\texttt{test\_auth} \& \texttt{test\_auth\_manager}):} Kiểm định cơ chế băm mật khẩu (Password Hashing) sử dụng thuật toán PBKDF2 kết hợp giá trị nhiễu (Salting), đồng thời kiểm tra sự nhất quán trong luồng ghi cơ sở dữ liệu nội bộ SQLite3.
    \item \textbf{Khối Kiểm thử Giao thức (\texttt{test\_protocol}):} Xác thực cơ chế Tuần tự hóa (Serialization) và Phân tích cú pháp (Parsing) của giao thức cấu trúc gói tin nhị phân. Các bài kiểm tra đặc biệt tập trung vào khả năng bắt ngoại lệ (Error Handling) khi đối mặt với các cấu trúc JSON bị biến dạng (Malformed JSON).
    \item \textbf{Khối Kiểm thử Giới hạn Dòng (\texttt{test\_ratelimit}):} Mô phỏng các xung lưu lượng mạng (Traffic Spikes) đột biến nhằm kiểm chứng năng lực phản xạ, từ chối dịch vụ của thuật toán Token Bucket tại biên giới mạng.
\end{itemize}

Bên cạnh các bài kiểm thử C++, dự án tích hợp hệ thống kiểm thử tự động hộp đen viết bằng ngôn ngữ Python (\texttt{integration\_test.py}). Kịch bản này tương tác trực tiếp qua luồng TCP (Raw Socket) nhằm kiểm chứng sự bền bỉ của kiến trúc đa luồng trước các trạng thái dị thường từ máy khách.

Bảng \ref{tab:test_suites} trình bày ma trận phân bổ các module kiểm thử tự động của hệ thống.

\begin{longtable}{
    >{\raggedright\arraybackslash}p{4.5cm}
    >{\centering\arraybackslash}p{3cm}
    >{\raggedright\arraybackslash}p{6cm}
}
    \caption{Ma trận phân bổ các module kiểm thử tự động} \label{tab:test_suites} \\
    \toprule
    \multicolumn{1}{c}{\textbf{Tên module kiểm thử}} & \multicolumn{1}{c}{\textbf{Ngôn ngữ lập trình}} & \multicolumn{1}{c}{\textbf{Mục tiêu kiểm định kỹ thuật}} \\
    \midrule
    \endfirsthead
    
    \toprule
    \multicolumn{1}{c}{\textbf{Tên module kiểm thử}} & \multicolumn{1}{c}{\textbf{Ngôn ngữ lập trình}} & \multicolumn{1}{c}{\textbf{Mục tiêu kiểm định kỹ thuật}} \\
    \midrule
    \endhead
    
    \bottomrule
    \endfoot
    
    \bottomrule
    \endlastfoot

    \texttt{test\_crypto} & C++ / Macro & Đảm bảo tính nhất quán của bộ mã hóa lõi \texttt{CryptoEngine}. \\
    \midrule
    
    \texttt{test\_protocol} & C++ / Macro & Đo lường và xác thực tính toàn vẹn gói tin (CRC32 Checksum) và cấu trúc Header. \\
    \midrule

    \texttt{test\_auth\_manager} & C++ / Macro & Xác thực khả năng chịu tải của Local DB và độ an toàn thuật toán PBKDF2. \\
    \midrule
    
    \texttt{test\_ratelimit} & C++ / Macro & Kiểm tra tính chính xác của cơ chế chống spam (Token Bucket Firewall). \\
    \midrule

    \texttt{integration\_test.py} & Python 3 & Kiểm thử tích hợp hộp đen (E2E), mô phỏng chu kỳ sống (Lifecycle) của Client. \\
\end{longtable}

[PLACEHOLDER\_HINH\_ANH\_KIEN\_TRUC\_KIEM\_THU]

\subsection[Đánh giá giới hạn chịu tải]{Đánh giá giới hạn chịu tải (Load Benchmark)}

\subsubsection[Cơ sở lý thuyết về đánh giá hiệu năng]{Cơ sở lý thuyết về đánh giá hiệu năng hệ thống mạng}

Đánh giá giới hạn chịu tải (Load Benchmark) là một hạng mục kiểm định cốt lõi đối với các kiến trúc phân tán và máy chủ truyền thông thời gian thực. Mục tiêu của quá trình này là đo lường hành vi của hệ thống dưới các mức độ áp lực dữ liệu (Data Pressure) khác nhau, từ đó phát hiện các điểm thắt cổ chai (Bottleneck) về I/O, khả năng tính toán của CPU hoặc giới hạn bộ nhớ.

Ba chỉ số định lượng (Metrics) quan trọng thường được sử dụng trong đánh giá hệ thống mạng bao gồm:
\begin{itemize}[label={-}]
    \item \textbf{Thông lượng (Throughput):} Đo lường số lượng gói tin hoặc luồng thông điệp mà hệ thống có thể xử lý thành công trong một đơn vị thời gian (thường là msg/s).
    \item \textbf{Độ trễ trung bình (Average Latency):} Đo lường thời gian trung bình hệ thống cần để hoàn thành chu trình từ khi tiếp nhận yêu cầu đến khi trả về kết quả (Round-Trip Time).
    \item \textbf{Độ trễ đuôi phân phối (Tail Latency - P99, P99.9):} Biểu thị thời gian phản hồi của 1\% hoặc 0.1\% lượng yêu cầu chậm nhất trong hệ thống. Đây là chỉ số khắt khe nhất nhằm đánh giá độ ổn định của kiến trúc đa luồng, đảm bảo luồng thực thi không bị phong tỏa (Blocked) đột biến.
\end{itemize}

\subsubsection{Kịch bản đo lường thực tế hệ thống}

Phép thử căng thẳng (Stress Test) dành cho mô hình bộ định tuyến sự kiện \texttt{epoll} được triển khai tự động thông qua mã nguồn \texttt{benchmark\_load.cpp}. Kịch bản này ứng dụng kỹ thuật Đa luồng (Multi-threading) để cấp phát đồng thời hàng loạt luồng ngầm độc lập, mỗi luồng vận hành như một máy khách (Client) kết nối trực tiếp thông qua Raw Socket đến cổng TCP \texttt{9000} của máy chủ.

Trong chu kỳ mô phỏng, mỗi máy khách sẽ gửi tuần tự gói tin \texttt{MSG\_CHAT\_SEND} dạng JSON định kỳ sau mỗi khoảng nghỉ 10 mili-giây. Kịch bản đánh giá được chia thành ba cấp độ (Scenarios) tăng dần:
\begin{itemize}[label={-}]
    \item \textbf{Baseline (Tiêu chuẩn):} Mô phỏng 10 máy khách hoạt động liên tục trong 30 giây.
    \item \textbf{Load Test (Chịu tải):} Tăng quy mô lên 50 máy khách hoạt động liên tục trong 60 giây.
    \item \textbf{Stress Test (Căng thẳng):} Tạo áp lực dội bom dữ liệu với 100 máy khách trong 60 giây.
\end{itemize}

Bảng \ref{tab:benchmark_performance} thống kê chi tiết các chỉ số hiệu suất đo lường được ở từng kịch bản.

\begin{longtable}{
    >{\raggedright\arraybackslash}p{3cm}
    >{\centering\arraybackslash}p{3cm}
    >{\centering\arraybackslash}p{3.5cm}
    >{\centering\arraybackslash}p{3cm}
    >{\centering\arraybackslash}p{3cm}
}
    \caption[Đo lường Thông lượng và Độ trễ của Server]{Đo lường Thông lượng và Độ trễ của Server ở các mức độ CCU} \label{tab:benchmark_performance} \\
    \toprule
    \multicolumn{1}{c}{\textbf{Kịch bản}} & \multicolumn{1}{c}{\textbf{Khách giả lập}} & \multicolumn{1}{c}{\textbf{Thông lượng}} & \multicolumn{1}{c}{\textbf{Độ trễ TB}} & \multicolumn{1}{c}{\textbf{Độ trễ P99}} \\
    \midrule
    \endfirsthead
    
    \toprule
    \multicolumn{1}{c}{\textbf{Kịch bản}} & \multicolumn{1}{c}{\textbf{Khách giả lập}} & \multicolumn{1}{c}{\textbf{Thông lượng}} & \multicolumn{1}{c}{\textbf{Độ trễ TB}} & \multicolumn{1}{c}{\textbf{Độ trễ P99}} \\
    \midrule
    \endhead
    
    \bottomrule
    \endfoot
    
    \bottomrule
    \endlastfoot

    Baseline & 10 CCU & $\sim$ 1,000 msg/s & 1.2 ms & 3.5 ms \\
    \midrule
    
    Load Test & 50 CCU & $\sim$ 4,800 msg/s & 2.8 ms & 8.1 ms \\
    \midrule

    Stress Test & 100 CCU & $\sim$ 9,200 msg/s & 5.4 ms & 15.3 ms \\
\end{longtable}

\subsubsection{Kết luận đánh giá hiệu năng}

Dưới áp lực tính toán phức hợp bao gồm việc tiếp nhận chuỗi kết nối I/O, thực thi giải mã luồng khối AES, phân tích và lọc dữ liệu (XSS Filtering) trước khi định tuyến phân phối lên đến \textbf{9,200 gói tin mỗi giây}, hệ thống vẫn duy trì chu kỳ sống trơn tru mà không ghi nhận bất kỳ hiện tượng tràn bộ nhớ hay treo cứng dịch vụ (Crash). 

Chỉ số độ trễ đuôi phân phối đối với 99\% lượng thông điệp (P99 Latency) vẫn duy trì ở ngưỡng cực đại $\sim$15.3 mili-giây, hoàn toàn đáp ứng được quy chuẩn kỹ thuật của một ứng dụng nhắn tin thời gian thực. Kết quả này xác thực tính đúng đắn của việc kết hợp kiến trúc cấp phát luồng dạng Thread-Pool cùng cơ chế điều phối phi đồng bộ \texttt{epoll\_wait} ở tầng nhân hệ điều hành Linux.

[PLACEHOLDER\_HINH\_ANH\_BIEU\_DO\_HIEU\_NANG]
